\hypertarget{appendix-12bbe185dbe5}{}\label{referencestart:appendix-12bbe185dbe5}
\section{Multivariate Gaussian distributions and whitening}\label{appendix-12bbe185dbe5-multivariate-gaussian-distributions-and-whitening}

\hypertarget{appendix-12bbe185dbe5-multivariate-gaussian-distributions-and-whitening}{}

\textbf{Curriculum depth:} Developed · \textbf{First seen:} Phase 10 · \textbf{Latest development:} Phase 10

\begin{allmightroledefinition}

\subsection{Definition and mechanism}\label{appendix-12bbe185dbe5-definition-and-mechanism}

\hypertarget{appendix-12bbe185dbe5-definition-and-mechanism}{}

A multivariate Gaussian is determined by a mean vector and a symmetric positive-semidefinite covariance matrix. When the covariance is positive definite, its density is an elliptical bell whose principal directions are covariance eigenvectors. Cholesky factorization maps independent standard-normal samples into correlated samples; whitening applies an inverse factor to express observations in unit-variance coordinates. Singular covariance instead confines the distribution to a lower-dimensional affine subspace and has no ordinary density with respect to full-dimensional volume. Covariance transforms as AΣAᵀ under a linear map and approximately as JΣJᵀ through a nonlinear map near a reference point. These equations support sensor fusion and projected uncertainty, but a Gaussian rendering kernel is not automatically a normalized probability density or epistemic uncertainty model.

\end{allmightroledefinition}

\subsection{Prerequisites}\label{appendix-12bbe185dbe5-prerequisites}

\hypertarget{appendix-12bbe185dbe5-prerequisites}{}

\begin{itemize}
\tightlist
\item
  \href{../05-Spatial-Worlds-Engine-State-and-Simulation-P01/main.pdf\#appendix-54c383df9992}{Expectation, covariance, and dependence (Spatial Worlds, Engine State, and Simulation, Part 1, Reference notes: Expectation, covariance, and dependence)}
\item
  \hyperref[appendix-50d2ea3c9b26]{Symmetric positive-definite matrices and Cholesky}
\item
  \href{../02-Mathematics-for-Visual-Computing-P02/main.pdf\#appendix-4815cd87ab8b}{Probability density and change of measure (Mathematics for Visual Computing, Part 2, Reference notes: Probability density and change of measure)}
\end{itemize}

\subsection{Used by}\label{appendix-12bbe185dbe5-used-by}

\hypertarget{appendix-12bbe185dbe5-used-by}{}

\begin{itemize}
\tightlist
\item
  \hyperref[appendix-f0437ea55a47]{Conditional Gaussians and Schur complements}
\item
  \hyperref[appendix-ad175ef3c5bf]{PCA and covariance eigensystems}
\end{itemize}

\subsection{Full teaching explanation}\label{appendix-12bbe185dbe5-full-teaching-explanation}

\hypertarget{appendix-12bbe185dbe5-full-teaching-explanation}{}

\begin{itemize}
\tightlist
\item
  \href{../02-Mathematics-for-Visual-Computing-P01/main.pdf\#02-foundations-ii-ch024-waiting-times-event-counts-and-covariance-geometry}{Phase 10 teaching section (Mathematics for Visual Computing, Part 1, Ch. 8)}
\item
  \href{../01-Bits-Programs-and-Reliable-Systems-P01/main.pdf\#appendix-cb5af353ef77}{Visual Recall/10 - Probability and Monte Carlo - Visual Recall.md (Bits, Programs, and Reliable Systems, Part 1, Reference notes: 10 - Probability and Monte Carlo - Visual Recall)} · visual recall
\end{itemize}

\subsection{Source coverage}\label{appendix-12bbe185dbe5-source-coverage}

\hypertarget{appendix-12bbe185dbe5-source-coverage}{}

Supplied lesson parts: 928 (T01 · M056), 929 (T01 · M056), 930 (T01 · M056).

Expanded direct teaching: \href{../02-Mathematics-for-Visual-Computing-P01/main.pdf\#02-foundations-ii-ch023-whiten-residuals-without-pretending-every-ellipse-has-the-same-probability}{Continuous Notes/Computing Foundations - From Bits to C++ and Systems APIs\#Whiten residuals without pretending every ellipse has the same probability (Mathematics for Visual Computing, Part 1, Ch. 7)}.
